\documentclass[11pt,a4paper]{article}
\usepackage[T1]{fontenc}
\usepackage{lmodern,amsmath,amssymb,amsthm,mathtools,booktabs,microtype}
\usepackage[margin=27mm]{geometry}
\usepackage[colorlinks=true,linkcolor=blue,citecolor=blue,urlcolor=blue]{hyperref}
\setlength{\emergencystretch}{2em}
\newtheorem{theorem}{Theorem}[section]
\newtheorem{lemma}[theorem]{Lemma}
\newtheorem{proposition}[theorem]{Proposition}
\newtheorem{corollary}[theorem]{Corollary}
\theoremstyle{remark}\newtheorem{remark}[theorem]{Remark}
\newcommand{\Q}{\mathbb Q}\newcommand{\Z}{\mathbb Z}
\newcommand{\R}{\mathbb R}\newcommand{\A}{\mathbb A}
\DeclareMathOperator{\N}{N}
\DeclareMathOperator{\Tr}{Tr}
\title{Fixed pole fibers and independent detector indices}
\author{}\date{16 November 2026}
\begin{document}\maketitle
\begin{abstract}
For a prepared pole input, one rational witness with nonzero ordinate at
a fixed conic parameter gives infinitely many witnesses at that same
parameter. The associated elliptic twist has odd valuation at 17, so its
rational torsion consists only of its two-torsion. Odd multiples then give
the required density. In the native detector sources, syntactic square
extraction produces positive scalar pole coefficients that are independent
of every auxiliary point and every finite-algebra character. Consequently
one common rational parameter with nonzero pole ordinates is sufficient
for all pole-stage guards. We also prove that, after one prepared tuple has
been chosen, the auxiliary cyclic indices can vary independently in
infinite return sets while preserving every prepared-input condition and
all point-stage guards. These statements reduce the shared-parameter
problem to its global rational existence assertion; they do not establish
that assertion or a smaller defining polynomial.
\end{abstract}

\section{Prepared inputs and fixed pole fibers}

A rational number $b$ is \emph{prepared} if its reduced denominator is a
square and
\begin{equation}\label{eq:preparation}
\begin{gathered}
0<b<17,\qquad
b\in\Q_2^{\times2}\cap\Q_3^{\times2}\cap\Q_{17}^{\times2},\\
v_2(b)\ge3,\qquad v_3(b)\ge1,\qquad v_{17}(b)\ge2.
\end{gathered}
\end{equation}
Write $b=dt^2$ with $d,t\in\Q^\times$. No squarefree hypothesis on
$d$ is imposed. For a rational parameter $\tau$ set
\begin{equation}\label{eq:normalized}
 D=17d,\qquad C=17(17-b),\qquad
 e=\frac{\tau^2-D}{\tau^2+3D},\qquad A=eC.
\end{equation}
The normalized second-pole equation of \cite{UnitNorm} is
\begin{equation}\label{eq:pole}
                  q^2=D(s^2-A)(s^2+3A).
\end{equation}
The normalized first conic $p^2=D(1-e)(1+3e)$ already has the
rational ordinate $p=-4D\tau/(\tau^2+3D)$. The original first
ordinate is $tp$. Neither ordinate is a retained coordinate or a
nonzero guard in the native source.

\begin{lemma}\label{lem:odd17}
For every prepared $b$ and every $\tau\in\Q$, the quantities in
\eqref{eq:normalized} are defined, $D>0$, $A\ne0$, and
\[
 v_{17}(e)=0,\qquad
 v_{17}(DA)=v_{17}(D)+2\quad\hbox{is odd}.
\]
In particular $DA$ is not a rational square.
\end{lemma}
\begin{proof}
Since $b=dt^2>0$, one has $d>0$ and $\tau^2+3D>0$.
The fact that $b$ is a square in $\Q_{17}$ gives even $v_{17}(d)$,
so $v_{17}(D)$ is odd. The valuations of $\tau^2$ and $D$ are
different, including when $\tau=0$. Consequently
\[
 v_{17}(\tau^2-D)=v_{17}(\tau^2+3D)
                 =\min\{2v_{17}(\tau),v_{17}(D)\},
\]
with the usual convention at zero. Thus $e$ is nonzero and has
valuation zero. Since $v_{17}(b)\ge2$, the number $17-b$ has
valuation one, and $v_{17}(C)=2$. The assertions follow.
\end{proof}

We record the torsion argument needed for fixed fibers. It applies to
negative as well as positive twists.

\begin{lemma}\label{lem:torsion}
If $n\in\Q^\times$ is not a rational square, then
\[
 E_n:\quad y^2=x(x-n)(x+3n)
\]
has rational torsion group $E_n[2](\Q)$.
\end{lemma}
\begin{proof}
There is full rational two-torsion. Mazur's theorem
\cite[Introduction, Theorem~(8), p.~35]{Mazur} leaves only
$\Z/2\Z\times\Z/(2m)\Z$, $1\le m\le4$.
It suffices to exclude orders three and four.
The three-division polynomial is
\[
                   3x^4+8nx^3-18n^2x^2-9n^4.
\]
A rational root would give a root of $3z^4+8z^3-18z^2-9$.
A root is integral at 5, because the leading coefficient is a unit
there; the values modulo 5 are $1,4,1,3,3$. Thus there is no
rational point of order three.

If a rational two-torsion point at a root $e_i$ is twice a rational
point, both differences $e_i-e_j$ and $e_i-e_k$ are rational squares;
this is the usual doubling criterion for a split cubic. At the root
zero, the differences $-n,3n$ have opposite signs. At the root
$-3n$, their ratio is $3/4$, which is not a square. At the root
$n$, the differences are $n,4n$, excluded by the hypothesis. There
is no point of order four. This also proves the torsion lemma used
in \cite{Density}.
\end{proof}

\begin{theorem}[Density in one fixed fiber]\label{thm:fixedfiber}
Fix a prepared $b=dt^2$ and a rational $\tau$. If
\eqref{eq:pole} has a rational solution with $q\ne0$, then the set
of its rational $s$ coordinates with $q\ne0$ is infinite. In
particular these witnesses can avoid the zeros of any finite collection
of nonzero polynomials in $\overline\Q[s]$.
\end{theorem}
\begin{proof}
Set $n=DA$. A solution gives a point on $E_n$ by
\begin{equation}\label{eq:ellipticmap}
                       (x,y)=(Ds^2,Dsq).
\end{equation}
Indeed substituting this pair in the elliptic equation gives
\eqref{eq:pole} multiplied by $D^2s^2$. The value $s=0$ is
impossible: the right side of \eqref{eq:pole} would be
$-3DA^2<0$. Thus $y\ne0$. Lemmas~\ref{lem:odd17} and
\ref{lem:torsion} show that this point has infinite order.

The square-class map
\[
 \alpha:E_n(\Q)\longrightarrow\Q^\times/\Q^{\times2},\qquad
 \alpha(O)=1,\quad \alpha((0,0))=-3n^2,\quad
 \alpha((x,y))=x\quad(x\ne0)
\]
is a homomorphism. One may verify this by the chord-and-tangent law:
the product of the three abscissas cut out by a line with nonzero
intercept is a square; the lines through $(0,0)$ give its displayed
exceptional value. Therefore every positive odd multiple of the point
\eqref{eq:ellipticmap} has abscissa in $D\Q^{\times2}$.
Its ordinate remains nonzero, and infinitely many different abscissas
occur because an abscissa determines at most two points. Write each
one as $Ds'^2$ and put $q'=y'/(Ds')$. This recovers infinitely many
rational $s'$ satisfying \eqref{eq:pole} with $q'\ne0$.
A nonzero one-variable polynomial has only finitely many roots.
\end{proof}

\begin{corollary}\label{cor:product}
For finitely many prepared inputs at the same fixed rational $\tau$,
suppose each second-pole equation has a rational solution with nonzero
ordinate. Their rational witness sets in the $s$ coordinates have
Zariski-dense product in the corresponding affine space. Any finite
collection of polynomial guards that remain nonzero polynomials after
fixing $\tau$ can be avoided simultaneously.
\end{corollary}
\begin{proof}
A product of infinite subsets of $\Q$ is Zariski dense: fix all
but one coordinate, apply the one-variable root bound, and repeat
for the coefficient polynomials. Apply this to the nonzero product
of the specialized guard polynomials.
\end{proof}

The specialization requirement is necessary. A polynomial that is
nonzero before $\tau$ is fixed can specialize to the zero polynomial.
The native pole guards satisfy the required stronger property, as shown
next.

\section{Scalar square classes in the native sources}

The native finite-algebra sources in \cite{UnitNorm} have the following
point roles and numbers of poles:
\begin{center}
\begin{tabular}{llll}\toprule
Source & Fixed root roles & Auxiliary roles & Poles\\\midrule
25 variables & $A,S,T,Z$ & $H,U_A,U_S,U_T,U_Z$ & 2\\
Explicit 27 variables & $A,S,T,Z$ & $H,U_A,U_S,U_T,U_Z$ & 3\\
Either 13-variable source & $A$ & $H,U_A$ & 2\\\bottomrule
\end{tabular}
\end{center}
Counts include the input. Let $B_0$ denote the tensor product of their
rank-four point algebras. Its base coordinates are the rational pairs
that specify the square roots of the point abscissas. Its finite
generators encode the possible point ordinates.

In the general and explicit sources the detector inputs are
\begin{equation}\label{eq:generalinput}
 b_i=(\rho R_i)
       \left(\frac{\N_{K/\Q}(\mathcal W_i)}{(\mathcal W_i)_0}\right)^2,
\end{equation}
whereas in either conditional source they are
\begin{equation}\label{eq:conditionalinput}
 b_i=F_i
       \left(\frac{\N_{K/\Q}(\mathcal W_i)}
                    {2\delta_i^4(\mathcal W_i)_0}\right)^2.
\end{equation}
Here $K=\Q(\sqrt2)$, and $(\mathcal W_i)_0$ is its rational
component. The coefficients $\rho R_i,F_i$ and the contact denominators
$\delta_i$ depend only on the input, root abscissas, and comparison
slope. The auxiliary points occur in $\mathcal W_i$.

\begin{proposition}\label{prop:scalar}
In all four native source circuits, syntactic square extraction writes
$b_i=d_i t_i^2$ with $d_i$ in the base scalar field. It depends only
on the input, root abscissas, comparison slope, and fixed source
parameters. It is independent of every auxiliary point and every
finite point-algebra character. After these earlier data have been
fixed, each $d_i$ is a fixed rational number. At a prepared actual
tuple, $d_i>0$.
\end{proposition}
\begin{proof}
The extraction in \cite{UnitNorm} flattens products and powers in
the numerator and denominator circuits. Every auxiliary or finite
ordinate dependence in \eqref{eq:generalinput} and
\eqref{eq:conditionalinput} occurs inside the displayed square, so
every such multiplicative atom has even exponent. It is assigned to
$t_i^2$. The odd exponents come entirely from $\rho R_i$ or $F_i$,
which are base scalars depending only on the stated earlier data.
Cancellation or a shared atom does not change this parity conclusion.
The argument is uniform in the allowed fixed exponent or shift of
each family. Finally, evaluating at the actual rational character
gives $b_i=d_it_i^2>0$ with $t_i\ne0$, whence $d_i>0$.
\end{proof}

In particular, the positivity conclusion holds at every point-algebra
character because the value of $d_i$ is identical at all of them. A
positive value at just one character would not suffice for this purpose.

\begin{theorem}[Pole guards at a common parameter]\label{thm:guards}
Fix an actual point-stage tuple of one of the four native sources for
which the detector inputs are prepared and all point-stage guards are
units in $B_0$. Suppose there is a common $\tau\in\Q$ for which
every actual second-pole equation has a rational solution with nonzero
ordinate. Then all native pole-stage guards can be satisfied without
changing $\tau$ or any earlier source data. This includes the norms
of all pole denominators, cleared radicands, and their syntactically
extracted factors.
\end{theorem}
\begin{proof}
By Proposition~\ref{prop:scalar}, each $d_i$ is a positive rational
scalar. Hence $\tau^2+51d_i>0$ at every character for every rational
$\tau$. All other denominators used to form the normalized pole
radicand are products of these conic denominators and the already
guarded fixed units.

For fixed $\tau$, the $i$th pole radicand in $B_0[s_i]$ is
\[
 17d_i\bigl(s_i^2-17e_i(17-b_i)\bigr)
       \bigl(s_i^2+51e_i(17-b_i)\bigr).
\]
Its leading coefficient is the unit $17d_i$. If $N=\dim_\Q B_0$,
the norm of this radicand is therefore a polynomial of degree $4N$
with leading coefficient $(17d_i)^N\ne0$. Clearing the already
invertible denominators preserves this nonvanishing statement.
The two displayed quadratic factors are monic, so their norms also
remain nonzero polynomials in $s_i$.

These are precisely the additional denominator and radicand types in
the native pole stage. Nonvanishing of a product at every character
implies nonvanishing of each syntactically extracted factor. In
particular, if a cleared radicand is written $S_i^2c_i$, its nonzero
norm makes both $S_i$ and $c_i$ units. Thus all required guard norms
are nonzero polynomials after the common $\tau$ is fixed.
Corollary~\ref{cor:product} now gives simultaneous avoidance by
independently varying the second-pole witnesses.
\end{proof}

The theorem concerns the pole stage only. Its hypothesis retains every
earlier point, difference, normalization, contact, and square-multiplier
guard. No condition on the first conic ordinate is added; in particular
$\tau=0$ is permitted whenever the stated rational second-pole witnesses
exist.

\section{Independent auxiliary indices after preparation}

Let
\[
 D_E=(3-\sqrt2)^2,\qquad E:\ y^2=x^3-D_E x,
\]
and retain the source premise $2E(K)=\Z P$, with the specified
infinite-order point $P$ of \cite{General,UnitNorm}. The retained point
chart has $x=h^2$, $y=hj$, with $h,j\in K$ and $j^2=h^4-D_E$.
The following assertion uses an already prepared point-stage tuple,
whose existence is the completeness conclusion of those sources.

\begin{samepage}
\begin{theorem}[Independent cyclic returns]\label{thm:indices}
Fix the input, root points, comparison slope, and an actual point-stage
tuple with all detector inputs prepared and every point-stage guard
nonzero at every character. Write its auxiliary points as
$R_{\ell,0}=[m_{\ell,0}]P$, $1\le\ell\le r$, where $r=5$ in
the general or explicit source and $r=2$ in a conditional source.
There are a positive integer $M$ and an infinite set
$\mathcal S\subset\Z_{>0}$ such that every independent choice
$(k_1,\ldots,k_r)\in\mathcal S^r$ gives auxiliary points
\begin{equation}\label{eq:independent}
             R_\ell=[m_{\ell,0}+2Mk_\ell]P
\end{equation}
with rational retained charts for which:
\begin{enumerate}
\item all earlier point-stage guards remain nonzero at every character;
\item the actual detector inputs remain prepared;
\item the root points, comparison scales, contacts, and extracted
      scalars $d_i$ are unchanged;
\item all exact difference relations are satisfied by recomputing
      $H-P$ and $U_i-A_i$.
\end{enumerate}
The choices can also preserve any prescribed finite collection of
finite-place congruences and real open conditions satisfied by the
starting tuple, and any fixed congruence conditions on its auxiliary
indices.
\end{theorem}
\end{samepage}
\begin{proof}
We first identify the conditions that must be preserved. The exact
source equations for auxiliary point data are their curve equations
and the rationally recovered differences $H-P$ and $U_i-A_i$.
The comparison functions, their ratios, and their scalar components
enter the detector products as functions; their vanishing is not an
additional source equation. In particular the relation
$U_i=[hn_i]P$ with $H=[h]P$ is a choice in the completeness proof,
not a relation imposed by the source. The comparison scales depend
only on the fixed input and root abscissas. These assertions follow
directly from the source formulas in \cite[Sections 3--5]{UnitNorm}.

Write each actual detector input as $b_i=c_i v_i^2$, with
$c_i=\rho R_i$ in \eqref{eq:generalinput} or $c_i=F_i$ in
\eqref{eq:conditionalinput}. The rational numbers $c_i\ne0$ are
now fixed. Let $S$ be a finite set of rational primes containing
$2,3,17$, every prime with odd valuation in some $c_i$, and all
primes appearing in the additional finite-place conditions. At
each prime in $S$, preparation persists in a sufficiently small
neighborhood of the starting tuple. Indeed local square classes and
the valuations of nonzero rational values are locally constant.
Outside $S$, every valuation of $b_i=c_iv_i^2$ is even. Any negative
valuation is consequently even, which is exactly what is needed
for the reduced denominator to remain a square. Thus preparation
requires only finitely many neighborhoods, together with the real
open inequalities $0<b_i<17$.

For all point-stage character guards, take the product of their
finitely many algebra norms after clearing already nonzero
denominators. It is a rational function of the retained base
coordinates with nonzero value at the starting tuple. Its
nonvanishing persists in a real neighborhood. This treats all
characters at once and requires no continuous choices of individual
ordinate signs.

We next give rational charts on which these continuity assertions
can be used. Since $R_{\ell,0}\in2E(K)$, choose
$Z_{\ell,0}\in E(K)$ with $2Z_{\ell,0}=R_{\ell,0}$. The
starting source guards exclude the identity and the points with
zero abscissa or ordinate. In particular $y(Z_{\ell,0})\ne0$.
For $Z=(x_Z,y_Z)$ the duplication identity is
\begin{equation}\label{eq:halfchart}
 x(2Z)=\left(\frac{x_Z^2+D_E}{2y_Z}\right)^2.
\end{equation}
\begin{samepage}
Choose $\varepsilon_\ell\in\{1,-1\}$ so that
\[
 h_\ell(Z)=\varepsilon_\ell\frac{x_Z^2+D_E}{2y_Z},\qquad
 j_\ell(Z)=\frac{y(2Z)}{h_\ell(Z)}
\]
agree with the original retained chart at $Z_{\ell,0}$.
Both are rational functions, regular in a neighborhood of that
half-point. Identity~\eqref{eq:halfchart} follows by substituting
$y_Z^2=x_Z^3-D_E x_Z$ in the usual duplication formula.
The signs change the chart, not the point $2Z$.
\end{samepage}

Compose all the required functions with these half-point charts.
Choose product neighborhoods of $(Z_{1,0},\ldots,Z_{r,0})$ at
the two real places and at \emph{every} place of $K$ over each
prime in $S$. By the preceding continuity arguments, they can be
chosen so that every tuple in them satisfies all the required
conditions. The exact differences are recomputed by their rational
formulas; their denominators stay nonzero on these neighborhoods.

At each of the finitely many nonarchimedean places, choose a
sufficiently small open subgroup of $E(K_v)$ whose translates
about all $Z_{\ell,0}$ lie in the chosen neighborhoods. Such a
subgroup has finite index, because $E(K_v)$ is compact. Choose
$M>0$ so that $[M]P$ belongs to all these subgroups, increasing
$M$ to be divisible by the desired index moduli. Then
$Z_{\ell,0}+[Mk]P$ belongs to the required finite neighborhoods
for every integer $k$ and every role $\ell$.

The product of the real elliptic groups is compact. There are
infinitely many positive $k$ for which $[Mk]P$ lies in any
specified neighborhood of its identity: take differences of widely
separated terms of a convergent subsequence of its positive orbit.
Choose the neighborhood small enough that its translate by each
$Z_{\ell,0}$ lies in the corresponding real neighborhood, and let
$\mathcal S$ be the resulting infinite set of return indices.
For independent $k_\ell\in\mathcal S$, put
$Z_\ell=Z_{\ell,0}+[Mk_\ell]P$ and $R_\ell=2Z_\ell$.
This gives \eqref{eq:independent} and rational charts with all
the required properties. Proposition~\ref{prop:scalar} shows
that the $d_i$ remain fixed. The half-points serve only to prove
existence of the charts; they are not added source coordinates.
\end{proof}

The theorem gives independent choices of cyclic indices, not arbitrary
rational point coordinates. It does not assert independent variation or
surjectivity of the resulting pair of detector values. Rational solubility
of the pole fibers is also not a condition whose preservation has been
proved by these neighborhoods.

\section{The remaining shared-parameter assertion}

For either two-pole source, the precise remaining sufficiency statement
is the following: for every integer input in its completeness range,
one can choose the fixed root and comparison data and an actual guarded
prepared auxiliary tuple such that some $\tau\in\Q$ makes both
equations
\[
 q_i^2=17d_i\left(s_i^2-17e_i(17-b_i)\right)
                   \left(s_i^2+51e_i(17-b_i)\right),\qquad
 e_i=\frac{\tau^2-17d_i}{\tau^2+51d_i},\qquad i=0,1,
\]
have rational solutions with $q_0q_1\ne0$.
Theorem~\ref{thm:indices} supplies prepared actual families with
independent auxiliary indices and fixed $d_i$.
Theorem~\ref{thm:guards} shows that no additional pole-guard
assertion is required once a common rational fiber is found.
The local theorem of \cite{Local} supplies points at every completion
for every rational $\tau$, but it does not supply rational points.
The global assertion just stated remains unproved.

If established with the retained arithmetic premises, it would permit
the identification of the two conic parameters, leaving the three
retained pole coordinates $\tau,s_0,s_1$. The unchanged point and
elimination coordinates would then give 24 total variables for the
general source and 12 for either conditional source. These numbers are
consequences of that unproved assertion, not bounds obtained here.

\newpage
\section{Exact source and identity checks}

The companion verifier is
\begin{center}\small
\texttt{computation/pole\_parameters/fixed\_fiber\_guards\_v42/check.py}.
\end{center}
It independently traverses the four native coefficient circuits and
checks the dependency sets of every extracted $d_i$. It also checks
$b_i=d_it_i^2$ by a separate multiplicative normal form, verifies
the half-point and pole-to-elliptic identities, and recovers exact
rational pole witnesses from the first four positive odd multiples
in six prepared examples with both signs of $DA$. Auxiliary checks
record the nonsquare-twist three-division obstruction and the
exceptional rational-square twist, which is not needed in
Lemma~\ref{lem:odd17}.

The four circuit runs use their default fixed parameters: exponent
$3$ for the general source, exponent $20483$ for the explicit source,
and shift $3$ for the two conditional sources. Uniformity in the
allowed fixed parameters follows from the proof of
Proposition~\ref{prop:scalar}, not from these finite runs. The verifier
does not establish the source arithmetic premises, infinite recurrence,
or simultaneous global rational realization.

\begin{thebibliography}{9}
\raggedright
\bibitem{General}
\emph{A defining polynomial in twenty-five variables:
combined scalar and comparison conditions},
project source, version 19,
\texttt{papers/few\_variables25/}\allowbreak
\texttt{few\_variables25\_merged\_20261110\_v19.tex}.
\bibitem{UnitNorm}
\emph{A single norm determinant for the 25- and 13-variable constructions},
project source, version 28,
\texttt{papers/finite\_algebra/}\allowbreak
\texttt{unit\_norm\_20261114\_v28.tex}.
\bibitem{Density}
\emph{Zariski density of the prepared pole witnesses},
project source, version 14,
\texttt{papers/pole\_parameters/}\allowbreak
\texttt{pole\_witness\_density\_20261109\_v14.tex}.
\bibitem{Local}
\emph{Local solubility at every shared pole parameter},
project source, version 41,
\texttt{papers/pole\_parameters/}\allowbreak
\texttt{shared\_parameter\_local\_solubility\_20261115\_v41.tex}.
\bibitem{Mazur}
B.~Mazur, \emph{Modular curves and the Eisenstein ideal},
Publications Math\'ematiques de l'IH\'ES \textbf{47} (1977), 33--186,
Introduction, Theorem~(8), p.~35,
\url{https://doi.org/10.1007/BF02684339}.
\end{thebibliography}
\end{document}
